\subsection{Weighted distance from a sum of multiplier ideals}
\label{subsec:psh-quantitative-ideals}

Multiplier ideals record local integrability. On a pseudoconvex domain, a weighted norm also measures how much of a holomorphic function remains after subtracting global sections of these ideals. The single-weight estimate can be applied successively, leaving the later weights in the auxiliary measure at each step.

Let $\Omega\subset\mathbb C^n$ be pseudoconvex, with $n\geq1$, and write $\mathrm{d}V$ for Euclidean volume. For a genuine psh weight $\rho$, put
\[
 \|f\|_\rho^2=\int_\Omega |f|^2\mathrm e^{-\rho}\,\mathrm{d}V.
\]
For a psh function $v$, $\Gamma(\Omega,\mathcal I(v))$ denotes the holomorphic functions whose germs belong to $\mathcal I(v)$ at every point. Membership requires local weighted integrability, without a finite global norm.

We first establish the two-summand estimate needed below. Its sharp constant comes from retaining the sum of the twisting weights in the $\bar\partial$ estimate.

\Needspace{10\baselineskip}
\begin{lemma}\label{lem:two-weight-l2-division}\index{Multiplier ideal!weighted division}
Let $\psi,\phi_1,\phi_2$ be genuine psh functions on $\Omega$, and suppose $S=\mathrm e^{\phi_1}+\mathrm e^{\phi_2}<1$. If $f$ is holomorphic and
\[
 J=\int_\Omega\frac{|f|^2\mathrm e^{-\psi}}{S^2}\,\mathrm{d}V<+\infty,
\]
then there are holomorphic $F_1,F_2$ with $F_1+F_2=f$ and
\begin{equation}\label{eq:two-weight-l2-division}
 \int_\Omega\bigl(|F_1|^2\mathrm e^{-\phi_1}
                  +|F_2|^2\mathrm e^{-\phi_2}\bigr)
                    \mathrm e^{-\psi}\,\mathrm{d}V\leq J.
\end{equation}
\end{lemma}
\begin{proof}
First suppose that the weights are smooth on a neighborhood of the closure of a bounded pseudoconvex domain, with $S<1$ there, and that $\Phi=\log S$ is strictly psh. Set $\chi=\phi_1+\phi_2+\psi$ and $A=(\Phi_{j\bar k})$. For a $(0,1)$-form $b$, write $|b|_A^2$ for its squared norm in the dual of this positive Hermitian matrix. We use
\begin{equation}\label{eq:negative-weight-dbar}
 \bar\partial u=b,\qquad
 \int |u|^2\mathrm e^{\Phi-\chi}\,\mathrm{d}V
 \leq\int (1-\mathrm e^\Phi)|b|_A^2
                         \mathrm e^{-\chi}\,\mathrm{d}V
\end{equation}
for $\bar\partial$-closed $b$ with finite right-hand side.

Here is the constant in this estimate. For $x<0$, define
\[
 T(x)=-\log(1-\mathrm e^x),\qquad
 \eta(x)=\frac{-x}{1-\mathrm e^x}-1,\qquad
 \lambda(x)=-\frac{\eta'(x)}{T'(x)}.
\]
The inequality $-x>1-\mathrm e^x$ gives $\eta>0$. Also
$\eta'= -[1-(1-x)\mathrm e^x]/(1-\mathrm e^x)^2<0$, whereas $T'>0$, so $\lambda>0$.
Direct differentiation gives
\[
 \eta T'-\eta'=1,\qquad
 \eta T''-\eta''-\frac{(\eta')^2}{\lambda}=0,\qquad
 \frac{\mathrm e^{-T}}{\eta+\lambda}=\mathrm e^x.
\]
Use the line metric $h=\mathrm e^{-\chi-T(\Phi)}$ and evaluate the twisting functions at $\Phi$. The twisted curvature is
\[
 \begin{aligned}
 B_0&=\eta\,\mathrm i\partial\bar\partial(\chi+T(\Phi))
       -\mathrm i\partial\bar\partial\eta
       -\mathrm i\lambda^{-1}\partial\eta\wedge\bar\partial\eta\\
    &=\eta\,\mathrm i\partial\bar\partial\chi
        +\mathrm i\partial\bar\partial\Phi
     \geq\mathrm i\partial\bar\partial\Phi.
 \end{aligned}
\]
Let $\Lambda$ denote contraction by the K\"ahler form and $\bar\partial_h^*$ the weighted adjoint. For a compactly supported line-valued $(n,1)$-form $v$, the twisted Bochner identity and weighted Cauchy inequality give
\begin{equation}\label{eq:sharp-twisted-apriori}
 \langle[B_0,\Lambda]v,v\rangle_h
 \leq\|\sqrt{\eta+\lambda}\,\bar\partial_h^*v\|_h^2
       +\|\sqrt\eta\,\bar\partial v\|_h^2.
\end{equation}
More explicitly, the cross term is bounded by
$\lambda|\bar\partial_h^*v|^2+
 \lambda^{-1}|(\bar\partial\eta)^*v|^2$;
the latter term is exactly the commutator contribution of
$\mathrm i\lambda^{-1}\partial\eta\wedge\bar\partial\eta$.
The terms of type $(n+1,1)$ vanish and the remaining connection norm has a favorable sign.

\Needspace{5\baselineskip}
The Hilbert space existence argument applied to \eqref{eq:sharp-twisted-apriori} therefore gives
\[
 \int\frac{|u|^2h}{\eta+\lambda}\,\mathrm{d}V
 \leq\int |b|_A^2h\,\mathrm{d}V.
\]
To see the coefficient, project a test form to $\ker\bar\partial$. On this kernel the second term of \eqref{eq:sharp-twisted-apriori} vanishes, and Cauchy inequality in the curvature norm bounds the functional $v\mapsto\langle b,v\rangle_h$ by $(\int |b|_A^2h)^{1/2}\|\sqrt{\eta+\lambda}\,\bar\partial_h^*v\|_h$. Riesz representation gives a solution with the displayed norm. The twisting functions are bounded on the present domain since $\Phi$ is smooth and bounded away from both $-\infty$ and $0$. The argument first uses a complete K\"ahler metric; approximation by the Euclidean metric plus a positive multiple of a complete metric gives the Euclidean statement. Under the identification with canonical-valued forms, the $(n,0)$ norm integral and the curvature-dual $(n,1)$ integral have these same coordinate expressions. The identities for $T,\eta,\lambda$ now give \eqref{eq:negative-weight-dbar}.

Put $v_i=\mathrm e^{\phi_i}/S$ and $\ell=(v_1f,v_2f)$. Equip $\mathbb C^2$ with the diagonal metric $(\mathrm e^{-\phi_1},\mathrm e^{-\phi_2})$. Then $\ell$ is orthogonal to the kernel line of the sum map, generated by $(1,-1)$, and $|\ell|^2=|f|^2/S$. Write $\delta=\phi_1-\phi_2$. The Hessian formula gives
\[
 \begin{aligned}
 A&=v_1(\phi_{1,j\bar k})+v_2(\phi_{2,j\bar k})
       +v_1v_2(\delta_j\overline{\delta_k}),\\
 \bar\partial(v_1f)&=fv_1v_2\bar\partial\delta.
 \end{aligned}
\]
Thus $v_1v_2|\bar\partial\delta|_A^2\leq1$ and, for $b=\bar\partial(v_1f)$,
\[
 |b|_A^2\mathrm e^{-\chi}
 \leq\frac{|f|^2\mathrm e^{-\psi}}{S^2}.
\]
Solve \eqref{eq:negative-weight-dbar} and set $F=\ell-u(1,-1)$. It is holomorphic and its components sum to $f$. The kernel norm is $\mathrm e^{-\phi_1}+\mathrm e^{-\phi_2}=S\mathrm e^{-\phi_1-\phi_2}$. Orthogonality consequently gives
\[
 \begin{aligned}
 \int |F|^2\mathrm e^{-\psi}
 &=\int\frac{|f|^2\mathrm e^{-\psi}}{S}
    +\int |u|^2\mathrm e^{\Phi-\chi}\\
 &\leq\int\frac{|f|^2\mathrm e^{-\psi}}{S^2}.
 \end{aligned}
\]

For general weights, exhaust $\Omega$ by bounded relatively compact pseudoconvex domains. Convolve each weight on a slightly larger neighborhood and add a sufficiently small positive multiple of $|z|^2$. These smooth psh weights dominate the original weights, converge to them almost everywhere, preserve $S<1$ on the chosen closure, and make $\Phi$ strictly psh. Their right-hand sides are bounded by $J$. On every compact subset the weights have common upper bounds, so \eqref{eq:two-weight-l2-division} gives uniform unweighted local $L^2$ bounds for both components. The holomorphic mean-value inequality and a diagonal extraction give holomorphic limits. Almost-everywhere convergence of the weights and Fatou's lemma yield \eqref{eq:two-weight-l2-division} on $\Omega$, while $F_1+F_2=f$ passes to the limit.
\end{proof}

\Needspace{10\baselineskip}
\begin{proposition}\label{prop:one-weight-ideal-distance}\index{Multiplier ideal!weighted distance}
Let $\psi$ be genuine psh, let $\phi<0$ be genuine psh, and let $p>1$. For holomorphic $f$ with $I=\|f\|_{\psi+\phi}^2<+\infty$, define
\[
 c=\inf\{\|g\|_\psi^2:g\in\mathcal O(\Omega),\quad
                   f-g\in\Gamma(\Omega,\mathcal I(p\phi+\psi))\}.
\]
Then $c\leq(1-1/p)I$.
\end{proposition}
\begin{proof}
For $t>0$, apply \cref{lem:two-weight-l2-division} with
\[
 \phi_1=\log\frac{t}{1+t}+p\phi,\qquad
 \phi_2=-\log(1+t).
\]
Their exponential sum is $(1+t\mathrm e^{p\phi})/(1+t)<1$. Its right-hand side is finite since $\|f\|_\psi^2\leq I$ and the sum is at least $1/(1+t)$. The first component belongs to $\Gamma(\Omega,\mathcal I(p\phi+\psi))$. The second is admissible for $c$, and the norm estimate gives
\[
 \frac{c}{1+t}\leq
   \int_\Omega\frac{|f|^2\mathrm e^{-\psi}}
                         {(1+t\mathrm e^{p\phi})^2}\,\mathrm{d}V.
\]
Set $t=x^p$ and integrate in $x>0$. With
$\theta_k=\int_0^\infty(1+x^p)^{-k}\,\mathrm{d}x$, Tonelli and the change of variable $x\mathrm e^\phi$ give $c\theta_1\leq I\theta_2$. The pole set is Lebesgue null. Since $p>1$, both integrals are finite, and integrating the derivative of $x/(1+x^p)$ gives $\theta_2=(1-1/p)\theta_1$. No attainment of the infimum is needed.
\end{proof}

\Needspace{18\baselineskip}
\begin{theorem}\label{thm:product-ideal-distance}\index{Multiplier ideal!sum of ideals}
Let $\psi$ and $\phi_1,\ldots,\phi_r<0$ be genuine psh functions on $\Omega$, with $r\geq1$, and choose $p_i>1$. For holomorphic $f$ with
\[
 I=\|f\|_{\psi+\sum_i\phi_i}^2<+\infty,
\]
put
\[
 c=\inf\left\{\|g\|_\psi^2:g\in\mathcal O(\Omega),\quad
 f-g\in\sum_{i=1}^r\Gamma(\Omega,\mathcal I(p_i\phi_i+\psi))\right\}.
\]
Then
\begin{equation}\label{eq:product-ideal-distance}
 c\leq\prod_{i=1}^r\left(1-\frac1{p_i}\right)I.
\end{equation}
For $r\leq n$, this constant is optimal for the stated class of data.
\end{theorem}
\begin{proof}
Set $\psi_i=\psi+\sum_{j>i}\phi_j$, $a_i=1-1/p_i$, and $f_0=f$. For arbitrary $\varepsilon_i>0$, \cref{prop:one-weight-ideal-distance} successively gives holomorphic $f_i$ such that
\[
 \begin{aligned}
 \|f_i\|_{\psi_i}^2&\leq a_i\|f_{i-1}\|_{\psi_i+\phi_i}^2+\varepsilon_i,\\
 f_{i-1}-f_i&\in\Gamma(\Omega,\mathcal I(p_i\phi_i+\psi_i)).
 \end{aligned}
\]
\Needspace{8\baselineskip}
Since $\psi_i\leq\psi$, these ideals are contained in $\mathcal I(p_i\phi_i+\psi)$. The sum of the remainders makes $f_r$ admissible for $c$. Also $\psi_i+\phi_i=\psi_{i-1}$ for $i>1$, so iteration gives
\[
 c\leq\left(\prod_i a_i\right)I
          +\sum_{i=1}^r\varepsilon_i\prod_{j>i}a_j.
\]
Let every $\varepsilon_i$ tend to zero.

\Needspace{9\baselineskip}
For sharpness when $r\leq n$, take $\Omega=\Delta^n$, $f=1$, $\psi=0$, and $\phi_i=p_i^{-1}\log|z_i|^2$ for $i\leq r$. By \cref{lem:smooth-center-multiplier-ideal}, $\mathcal I(p_i\phi_i)=(z_i)$. Every admissible $g$ equals one on $Y=\{z_1=\cdots=z_r=0\}$, hence $g(0)=1$. Orthogonality of polydisc monomials gives $\|g\|_0^2\geq\pi^n$, with equality for $g=1$. Direct radial integration gives $I=\pi^n\prod_i p_i/(p_i-1)$. Equality follows in \eqref{eq:product-ideal-distance}.
\end{proof}

The congruence in this theorem uses a sum of global sections. A condition at one germ alone does not supply these successive global remainders. The estimate holds for any number of weights; the optimality assertion uses distinct coordinate directions and applies only when $r\leq n$.
